\addcontentsline{toc}{section}{\normalsize{\bengalifont সারসংক্ষেপ}}
\begin{center}
{\bengalifont\Large সারসংক্ষেপ}\\[24pt]
{\bengalifont\bfseries সাবচ্যাট ট্রিজ: এলএলএম-এ মাল্টি-থ্রেডেড সংলাপ ও প্রসঙ্গ বিচ্ছিন্নকরণের জন্য একটি স্কেলযোগ্য স্তরবিন্যস্ত স্থাপত্য}
\end{center}
{\bengalifont
\justifying
বর্তমান বৃহৎ ভাষা মডেলভিত্তিক (LLM) চ্যাট ইন্টারফেসগুলো ব্যবহারকারীর সঙ্গে যোগাযোগের জন্য রৈখিক কথোপকথন-ইতিহাস ব্যবহার করে। রৈখিক কথোপকথনে প্রতিটি বিষয়, উপবিষয় এবং অনুসরণমূলক প্রশ্ন একই কালানুক্রমিক ইতিহাসে যুক্ত হয়। স্বল্পমেয়াদি ও একক-বিষয়ক কথোপকথনের ক্ষেত্রে এই নকশা কার্যকর হলেও বাস্তবসম্মত দীর্ঘমেয়াদি কথোপকথনে এটি দুর্বল হয়ে পড়ে, বিশেষত যখন ব্যবহারকারী কাজ বা বিষয় পরিবর্তন করেন, আগের কোনো বিষয়ে ফিরে যান বা স্তরবিন্যস্ত অনুসরণমূলক প্রশ্ন করেন। মানুষের কথোপকথন স্বাভাবিকভাবেই স্তরবিন্যস্ত; কিন্তু আধুনিক কৃত্রিম বুদ্ধিমত্তাভিত্তিক ইন্টারফেসে কথোপকথনের প্রবাহ নিয়ন্ত্রণের পর্যাপ্ত ব্যবস্থা নেই। ফলে একটি বিষয়ের সঙ্গে সম্পর্কহীন প্রসঙ্গ অন্য বিষয়ে প্রবেশ করতে পারে এবং পুরোনো প্রসঙ্গ সক্রিয় ইতিহাস থেকে বাদ পড়তে বা বিপুল তথ্যের মধ্যে চাপা পড়তে পারে। বর্তমান গবেষণাগুলো সাধারণত কনটেক্সটের দৈর্ঘ্য ও মডেলের জটিলতা বাড়িয়ে এই সমস্যা সমাধানের চেষ্টা করে। তবে শুধু এগুলো বাড়ালেই চ্যাট ইন্টারফেসে মানব কথোপকথনের স্বাভাবিক প্রবাহ বজায় থাকে না। আমরা সমস্যাটিকে একটি নকশাগত সীমাবদ্ধতা হিসেবে বিবেচনা করে কথোপকথন ব্যবস্থাপনার জন্য ব্যবহারকারীনির্ভর স্তরবিন্যস্ত শাখা-ভিত্তিক স্থাপত্য, Subchat Trees, প্রস্তাব করেছি। একটি সমতল রৈখিক কথোপকথন হিসেবে বিবেচনা না করে Subchat Trees সংলাপকে প্যারেন্ট ও চাইল্ড সাবচ্যাটের সমন্বয়ে গঠিত একটি ট্রি হিসেবে উপস্থাপন করে, যেখানে প্রতিটি নোড নিজস্ব স্থানীয় বাফার, ধারাবাহিক সারসংক্ষেপ এবং একটি অভিন্ন ভেক্টর-ইনডেক্সড দীর্ঘমেয়াদি মেমরি বজায় রাখে। ব্যবহারকারী প্যারেন্ট কথোপকথনের নির্বাচিত অংশ থেকে একটি নির্দিষ্ট বিষয়ে কেন্দ্রীভূত সাবচ্যাট তৈরি করতে পারেন, সেই শাখায় বিস্তারিত পার্শ্ব-আলোচনা চালিয়ে যেতে পারেন এবং পরে সমপর্যায়ের বা পূর্বসূরি প্রসঙ্গে অনাকাঙ্ক্ষিত তথ্য না মিশিয়ে মূল কথোপকথনে ফিরে আসতে পারেন। এই বাহ্যিক স্থাপত্যটি মডেল-নিরপেক্ষ এবং ফাইন-টিউনিং ছাড়াই LLM-এর কর্মদক্ষতা উন্নত করতে পারে, কারণ ইনফারেন্সের সময় মডেলটি সংশ্লিষ্ট শাখার তুলনামূলকভাবে পরিচ্ছন্ন ও প্রাসঙ্গিক কনটেক্সট পায়। প্রস্তাবিত পদ্ধতিটি পাঁচটি LMSYS-Chat-1M কথোপকথন থেকে তৈরি নিয়ন্ত্রিত 353টি টার্নের Variable-Block Randomized Interleaving (VBRI) স্ট্রেস-টেস্ট বেঞ্চমার্কে মূল্যায়ন করা হয়েছে, যেখানে 59টি বিষয়-লেবেল এবং 43টি বিষয়-স্মরণ যাচাইকারী প্রোব অন্তর্ভুক্ত রয়েছে। Qwen2.5 7B, 14B ও 32B মডেল এবং $C \in \{5,10,20,40\}$ বাফার আকারের প্রতিটি পরীক্ষিত বিন্যাসে Subchat Trees রৈখিক বেসলাইনের তুলনায় কনটেক্সট-আইসোলেশন F1 উন্নত করেছে। সর্বোচ্চ কর্মদক্ষতা পাওয়া গেছে $C=20$-এ Qwen2.5 32B মডেলে, যেখানে F1 স্কোর ছিল 85.2\% এবং কনটেক্সট দূষণের হার ছিল 16.2\%। স্মরণ যাচাইকারী প্রোবগুলোও বিষয়ভিত্তিক তথ্য ধরে রাখার উন্নত সক্ষমতা দেখিয়েছে। সব পরীক্ষিত বিন্যাসে বেসলাইনের তুলনায় গড় ROUGE-1 বৃদ্ধি পেয়েছে; এর মধ্যে $C=20$-এ 14B মডেলের স্কোর 0.2476 থেকে 0.4187 এবং 32B মডেলের স্কোর 0.3338 থেকে 0.4283 হয়েছে। একই সঙ্গে, সব পরীক্ষিত বিন্যাসে প্রতিটি সঠিক উত্তরের জন্য প্রয়োজনীয় ইনপুট টোকেনের সংখ্যাও কমেছে। ফলাফলগুলো নির্দেশ করে যে জটিল কথোপকথন কার্যকরভাবে পরিচালনার জন্য শুধু দীর্ঘ কনটেক্সট উইন্ডো যথেষ্ট নয়; এর পাশাপাশি সুস্পষ্ট বাহ্যিক কাঠামোও প্রয়োজন।
\par
}
